\section{Planteamiento del Problema}
    Para un conjunto de imágenes se deberá comparar:
    \begin{enumerate}
        \item Métodos de detección
        \item Ventajas y desventajas de los diferentes métodos
    \end{enumerate}
    Así como comparar el resultado de aplicar los siguientes 
    tipos de suavizados sobre la imagen (para cada método de 
    detección), analizando el efecto sobre los bordes:
    \begin{itemize}
        \item Sin suavizado.
        \item Con suavizado gaussiano con un kernel de $3\times3$ píxeles.
        \item Con suavizado gaussiano con un kernel de $7\times7$ píxeles.
    \end{itemize}

    Adicionalmente, se deberá incluir la interpretación de las 
    siguientes expresiones:
    \begin{itemize}
        \item $G_x = \frac{\partial I(x, y)}{\partial x}$
        \item $G_y = \frac{\partial I(x, y)}{\partial y}$
        \item $G = \sqrt{G_x^2+G_y^2}$
        \item $\theta = \arctan(\frac{G_y}{G_x})$
    \end{itemize}

    Por último, se visualizará la orientación de los bordes colocando 
    una flecha con inicio en el pixel actual, de longitud conveniente, 
    que apunte en la dirección señalada por el ángulo calculado $\theta$.

    \subsection{Restricciones}
        Al menos Sobel y Prewitt deberán implementarse manualmente con 
        NumPy, sin utilizar una función de alto nivel que haga el 
        procedimiento.